\documentclass[a4paper,10pt]{article}
\usepackage[left=3cm,top=3cm,right=2.5cm,bottom=4cm]{geometry}
\usepackage[utf8]{inputenc}
\usepackage[activeacute,spanish]{babel}
\usepackage{multirow}	% Para poder unir filas en las tablas
\usepackage{colortbl}	% Para colorear tablas
\usepackage{latexsym,amssymb,oldlfont,amsmath}
\usepackage{graphicx,graphics}
\usepackage{multicol}
\usepackage{color}
\usepackage{listings}
\usepackage{lastpage}
\usepackage{enumitem}
\usepackage{tkz-euclide}
\usetkzobj{all}


%\usepackage{fontspec}
%\setmainfont{Times New Roman}
\newcommand{\R}{{\mathbb R}}
\newcommand{\si}{{\mathcal S}^n}
\renewcommand{\theequation}{\thesection.\arabic{equation}} % Cambia formato salida numeracion equation
\makeatletter\@addtoreset {equation}{section} \makeatother % Tambien se usa
\renewcommand{\baselinestretch}{1.3}
\usepackage{amssymb}
\usepackage{hyperref}
\definecolor{midnightblue}{rgb}{0.1,0.1,0.44}


\newcommand{\dis}{\displaystyle}
\renewcommand{\theenumii}{\arabic{enumii}}

\def\Z{\mathbb{Z}}
\def\N{\mathbb{N}}
\def\R{\mathbb{R}}
\def\M{\mathcal{M}_n(\mathbb{R})}
\def\sin{\operatorname{sen}}
\def\e{\operatorname{e}}
\DeclareMathOperator{\arcsec}{arcsec}
\DeclareMathOperator{\arccot}{arccot}
\DeclareMathOperator{\arccsc}{arccsc}



\usepackage{etoolbox,fancyhdr,xcolor}
\fancyhf{} % sets both header and footer to nothing
\renewcommand{\headrulewidth}{0pt}
%\renewcommand{\footrulewidth}{0.5pt}
\pagestyle{fancy}
\headheight=60pt %para cambiar el tamaño del encabezadox
\fancyhead[L] %la "L" indica a la izquierda
{
\begin{minipage}{6cm}
\includegraphics[width=0.9\textwidth]{logo_udp.eps}
\end{minipage}
}

%\fancyhead[C]
%{
%\textbf{Cálculo II - Guía 1}
%}


%\fancyhead[R] %la "L" indica a la izquierda
%{

%}
%\begin{minipage}{4cm}
%\small %tamaño de letra pequeño
%{
%\textsf
%{
%Universidad de Chile\\
%...}}
%\end{minipage}}



%\fancyfoot[C]{}
%{
%\begin{minipage}{10cm}
%\color{midnightblue}
%\begin{center}Casa Matriz: Monjitas 527 Piso 15 - Santiago de Chile\\
%Fono (02)24367255 \end{center}
%\end{minipage}
%}
\fancyfoot[C]{}
%Página \thepage\ de \pageref{LastPage}}% Páginas Pares numeradas a la derecha
%\fancyfoot[LO]{\thepage} % Páginas Impares numeradas a la izquierda

\title{C2_Ay2_pauta}

\begin{document}

\begin{center}
Cálculo II\\
Ayudantía Nº2 - Pauta\\
Primer Semestre 2017
\end{center}


\begin{enumerate}


\item Calcule las siguientes integrales de funciones
trigonom\'{e}tricas.
\begin{enumerate}[label=\emph{\alph*)}]
\item $\dis \int x\sin^3(x^2)\, dx$\\

Utilizamos la siguiente sustitución:
\begin{align*}
 u  &=  x^2  &    \\
 du &= 2x  \, dx     &   x \, dx &= \frac{du}{2}
\end{align*}
\dotfill [1 punto]


Obtenemos lo siguiente:



\begin{align*}
 \int \frac{\sin^3(u)}{2} \, du  &=  \frac{1}{2} \int \sin(u)(1-\cos^2(u)) \, du\\
							     &=  \frac{1}{2} \int \sin(u) \, du - \frac{1}{2} \int \cos^2(u)\sin(u) \, du  \\
                                 &=  - \frac{1}{2}  \cos(u)  + \frac{1}{2} \frac{\cos^3(u)}{3} + C   \\
                                 &=  - \frac{1}{2}  \cos(x^2)  + \frac{1}{2} \frac{\cos^3(x^2)}{3} + C   \\
\end{align*}
\dotfill [1 punto]

\newpage
\item $\dis \int \frac{1-\sin x}{\cos x}\, dx$\\



\begin{align*}
 \dis \int \frac{1-\sin x}{\cos x}\, dx  &=    \int \frac{(1-\sin x)}{\cos x}\frac{(1+\sin x)}{(1+\sin x)}  \, dx\\
 										 &=    \int \frac{(1-\sin^2 x)}{\cos x(1+\sin x)}  \, dx\\
                                         &=    \int \frac{\cos^2 x}{\cos x(1+\sin x)}  \, dx\\
                                         &=    \int \frac{\cos x}{(1+\sin x)}  \, dx\\
\end{align*}


Utilizamos la siguiente sustitución:
\begin{align*}
u  &=  1 + \sin(x)      \\
du &= \cos(x)  \, dx
\end{align*}
\dotfill [1 punto]

Obtenemos lo siguiente:



\begin{align*}
 \dis \int \frac{1-\sin x}{\cos x}\, dx  &=  \int \frac{1}{u}  \, du\\
							     &=  \ln (u) + C \\
                                 &=  \ln (1+ \sin(x)) + C
\end{align*}

\end{enumerate}
\dotfill [1 punto]
\newpage

\item Calcule las siguientes integrales de funciones
trigonom\'{e}tricas.
\begin{enumerate}[label=\emph{\alph*)}]
\item $\dis \int \cot^5 x\sin^2 x\, dx$
\begin{align*}
 \dis \int \cot^5 x\sin^2 x\, dx  &=    \dis \int \frac{\cos^5 x}{\sin^5 x}\sin^2 x\, dx\\
 										 &=    \dis \int \frac{\cos^5 x}{\sin^3 x}\, dx\\
                                         &=    \dis \int \frac{\cos x (1-\sin^2 x)^2 }{\sin^3 x}\, dx
\end{align*}


Utilizamos la siguiente sustitución:
\begin{align*}
u  &= \sin(x)      \\
du &= \cos(x)  \, dx
\end{align*}
\dotfill [1 punto]

Obtenemos lo siguiente:

\begin{align*}
 \dis \int \frac{ (1-\sin^2 x)^2 }{\sin^3 x}\cos x \, dx  &=    \dis \int \frac{(1-u^2)^2 }{u^3}\, du\\
 										 &=    \dis \int \frac{(1-2u^2 + u^4) }{u^3}\, du\\
                                         &=    \dis \int \frac{1}{u^3 }\, dx - 2 \int \frac{u^2}{u^3 }\, dx + \int \frac{u^4}{u^3 }\, dx\\
                                         &=    \dis \int \frac{1}{u^3 }\, dx - 2 \int \frac{1}{u }\, dx + \int u\, dx\\
                                         &=    \dis -\frac{1}{2u^2 } - 2 \ln(u) +  u^2 + C\\
                                         &=    \dis -\frac{1}{2(\sin(x))^2 } - 2 \ln(\sin(x)) +  (\sin(x))^2 + C
\end{align*}
\dotfill [1 punto]

\newpage
\item $\dis \int \cot^4 x\csc^4 x\, dx$
\begin{align*}
 \dis \int \cot^4 x\csc^4 x\, dx  &=    \dis \int \cot^4 x\csc^2 x(1+\cot^2 x)\, dx\\
 										 &=    \dis \int \cot^4 x\csc^2 x\, dx +\dis \int \cot^6 x\csc^2 x\, dx
\end{align*}
Utilizamos la siguiente sustitución:
\begin{align*}
u  &= \cot(x)    &  \\
du &= -\csc^2(x)  \, dx  & \csc^2 (x) = -du
\end{align*}
\dotfill [1 punto]

Obtenemos lo siguiente:
\begin{align*}
 \dis \int \cot^4 x\csc^2 x\, dx +\dis \int \cot^6 x\csc^2 x\, dx  &=    -\dis \int u^4 \, du -\dis \int u^6 \, du\\
 										 &=    -\dis \frac{u^5}{5} -\dis \frac{u^7}{7} + C\\
                                          &=    -\dis \frac{(\cot (x))^5}{5} -\dis \frac{(\cot (x))^7}{7} + C
\end{align*}
\dotfill [1 punto]

\end{enumerate}


\newpage
\item Calcule las siguientes integrales de funciones
trigonom\'{e}tricas.
\begin{enumerate}[label=\emph{\alph*)}]
\item $\dis \int \sin^3 x\sqrt{\cos x}\, dx$
\begin{align*}
 \dis \int \sin^3 x\sqrt{\cos x}\, dx  &=    \dis \int (1-\cos^2 x)\sqrt{\cos x}\, \sin x \, dx
\end{align*}
Utilizamos la siguiente sustitución:
\begin{align*}
u  &= \cos(x)      \\
du &= -\sen(x)  \, dx
\end{align*}
\dotfill [1 punto]

Obtenemos lo siguiente:
\begin{align*}
 \dis \int (1-\cos^2 x)\sqrt{\cos x}\, \sin x \, dx  &=    -\dis \int (1-u^2)\sqrt{u} \, du\\
 													 &=    \dis -\int u^{\frac{1}{2}} \, du + \int u^{\frac{5}{2}} \, du\\
                                                     &=    \dis -\frac{u^{\frac{3}{2}}}{\frac{3}{2}} + \frac{u^{\frac{7}{2}}}{\frac{7}{2}} + C\\
                                                     &=    \dis -\frac{2}{3}u^{\frac{3}{2}} + \frac{2}{7}u^{\frac{7}{2}} + C\\
                                                     &=    \dis -\frac{2}{3}(\cos x)^{\frac{3}{2}} + \frac{2}{7}(\cos x)^{\frac{7}{2}} + C\\
\end{align*}
\dotfill [1 punto]


\newpage
\item $\dis \int \sin(5x) \sin(2x)\, dx$\\

Utilizando la siguiente identidad:
\begin{align*}
 \dis \sin A \cos B  &=    \frac{\sin (A+B) + \cos (A-B)}{2}
\end{align*}

reescribimos el integrando como $\dis \sin(5x)\cos(2x) = \frac{\sin(7x)+\cos(3x)}{2}$

\dotfill [1 punto]

Así, Oobtenemos lo siguiente:
\begin{align*}
 \dis \int \sin(5x) \sin(2x)\, dx  &=    \dis \int \frac{\sin (7x) + \cos (3x)}{2} \, dx\\
 								   &=    \dis \int \frac{\sin (7x)}{2} \, dx + \int \frac{\cos (3x)}{2} \, dx\\
                                   &=    \dis \frac{-\cos (7x)}{14} + \frac{\sin (3x)}{6} + C
\end{align*}
\dotfill [1 punto]




\end{enumerate}


\newpage
\item Calcule las siguientes integrales usando sustituci\'{o}n
trigonom\'{e}trica.
\begin{enumerate}[label=\emph{\alph*)}]
\item $\dis \int \frac{x^2}{\sqrt{5-x^2}}\, dx$\\

Utilizamos la siguiente sustitución trigonométrica:\\

\begin{minipage}{.3\textwidth}
\raggedright
\begin{align*}
x  &= \sqrt{5}\cos \theta      \\
dx &= -\sqrt{5}\sin \theta  \, d\theta
\end{align*}
\end{minipage}% <---------------- Note the use of "%"
\begin{minipage}{.3\textwidth}
\raggedleft
\begin{tikzpicture}[thick]
\coordinate (O) at (0,0);
\coordinate (A) at (4,0);
\coordinate (B) at (0,2);
\draw (O)--(A)--(B)--cycle;

\tkzLabelSegment[below=2pt](O,A){$x$}
\tkzLabelSegment[left=2pt](O,B){$\sqrt{5-x^2}$}
\tkzLabelSegment[above right=2pt](A,B){$\sqrt{5}$}

\tkzMarkAngle[fill= orange,size=0.8cm,%
opacity=.4](B,A,O)
\tkzLabelAngle[pos = 0.6](B,A,O){$\theta$}

\end{tikzpicture}
\end{minipage}


Obtenemos lo siguiente:
\begin{align*}
 \dis \int \frac{x^2}{\sqrt{5-x^2}}\, dx  &=   - \dis \int \frac{(\sqrt{5}\cos \theta)^2}{\sqrt{5}\sin \theta}\sqrt{5}\sin \theta  \, d\theta\\
 										  &=  -  \dis \int \cos^2 \theta  \, d\theta
\end{align*}
\dotfill [1 punto]

Utilizando la siguiente identidad:
\begin{align*}
 \dis \cos (2\theta)  &=    2\cos^2 \theta - 1\\
 \dis \cos^2 \theta  &=    \frac{\cos (2\theta)+1}{2}
\end{align*}
Entonces obtenemos lo siguiente:
\begin{align*}
 \dis \int \cos^2 \theta  \, d\theta  &=   - \dis \int \frac{\cos (2\theta)+1}{2}  \, d\theta\\
 									  &=  -  \dis \int \frac{\cos (2\theta)}{2}  \, d\theta - \int \frac{1}{2} \, d\theta\\
                                      &=    \dis -\frac{\sin (2\theta)}{4}  \, d\theta -  \frac{\theta}{2} + C\\
                                      &=    \dis -\frac{2\sin \theta \cos \theta}{4}  \, d\theta -  \frac{\theta}{2} + C\\
                                      &=    \dis - \frac{1}{2} (\frac{\sqrt{5-x^2}}{\sqrt{5}}\frac{x}{\sqrt{5}} + \arccos (\frac{x}{\sqrt{5}}))   + C\\
                                      &=    \dis - \frac{x\sqrt{5-x^2}}{10} - \frac{\arccos (\frac{x}{\sqrt{5}}) }{2}   + C
\end{align*}
\dotfill [1 punto]


\newpage
\item $\dis \int \sqrt{2x-x^2}\, dx$\\

\begin{align*}
 \dis \int \sqrt{2x-x^2}\, dx  &=    \dis \int \sqrt{1 - (x-1)^2}\, dx
\end{align*}

Utilizamos la siguiente sustitución trigonométrica:\\

\begin{minipage}{.3\textwidth}
\raggedright
\begin{align*}
x -1  &= \cos \theta      \\
dx &= -\sin \theta  \, d\theta
\end{align*}
\end{minipage}% <---------------- Note the use of "%"
\begin{minipage}{.3\textwidth}
\raggedleft
\begin{tikzpicture}[thick]
\coordinate (O) at (0,0);
\coordinate (A) at (4,0);
\coordinate (B) at (0,2);
\draw (O)--(A)--(B)--cycle;

\tkzLabelSegment[below=2pt](O,A){$x-1$}
\tkzLabelSegment[left=2pt](O,B){$\sqrt{1- (x-1)^2}$}
\tkzLabelSegment[above right=2pt](A,B){$1$}

\tkzMarkAngle[fill= orange,size=0.8cm,%
opacity=.4](B,A,O)
\tkzLabelAngle[pos = 0.6](B,A,O){$\theta$}

\end{tikzpicture}
\end{minipage}

Obtenemos lo siguiente:
\begin{align*}
 \dis \int \sqrt{1 - (x-1)^2}\, dx  &=    \dis \int \sin \theta (-\sin \theta)  \, d\theta\\
 									&=	 -\dis \int \sin^2 \theta  \, d\theta\\
\end{align*}
\dotfill [1 punto]

Utilizando la siguiente identidad:
\begin{align*}
 \dis \cos (2\theta)  &=    1 - 2\sin^2 \theta \\
 \dis \sin^2 \theta  &=    \frac{\cos (2\theta)-1}{2}
\end{align*}
Entonces obtenemos lo siguiente:
\begin{align*}
 -\dis \int \sin^2 \theta  \, d\theta  &=    -\dis \int \frac{\cos (2\theta)-1}{2}  \, d\theta \\
 									   &=	 -\dis \int \frac{\cos (2\theta)}{2}  \, d\theta + \int \frac{1}{2} \, d\theta\\
                                       &=	 -\dis \ \frac{\sin (2\theta)}{4}  +  \frac{\theta}{2} + C\\
                                       &=	 -\dis \ \frac{2\sin \theta \cos \theta}{4}  +  \frac{\theta}{2} + C\\
                                       &=    \dis - \frac{(x-1)\sqrt{1-(x-1)^2}}{2} + \frac{\arccos (x-1) }{2}   + C
\end{align*}

\end{enumerate}
\dotfill [1 punto]

\newpage
\item Calcule las siguientes integrales usando sustituci\'{o}n
trigonom\'{e}trica.
\begin{enumerate}[label=\emph{\alph*)}]
\item $\dis \int \frac{1}{x^2\sqrt{16x^2-9}}\, dx$\\

Utilizamos la siguiente sustitución trigonométrica:\\

\begin{minipage}{.3\textwidth}
\raggedright
\begin{align*}
x  &= \frac{3}{4} \sec \theta      \\
dx &= \frac{3}{4} \sec \theta \tan \theta  \, d\theta
\end{align*}
\end{minipage}% <---------------- Note the use of "%"
\begin{minipage}{.3\textwidth}
\raggedleft
\begin{tikzpicture}[thick]
\coordinate (O) at (0,0);
\coordinate (A) at (4,0);
\coordinate (B) at (0,2);
\draw (O)--(A)--(B)--cycle;

\tkzLabelSegment[below=2pt](O,A){$3$}
\tkzLabelSegment[left=2pt](O,B){$\sqrt{16x^2- 9}$}
\tkzLabelSegment[above right=2pt](A,B){$4x$}

\tkzMarkAngle[fill= orange,size=0.8cm,%
opacity=.4](B,A,O)
\tkzLabelAngle[pos = 0.6](B,A,O){$\theta$}

\end{tikzpicture}
\end{minipage}

\dotfill [1 punto]

Obtenemos lo siguiente:
\begin{align*}
 \dis \int \frac{1}{x^2\sqrt{16x^2-9}}\, dx  &=    \dis \int \frac{1}{(\frac{3}{4}\sec \theta)^2\, (3\tan \theta)} \frac{3}{4}\sec \theta \tan \theta \, d\theta\\
 									&=    \dis \frac{4}{9} \int \cos \theta \, d\theta\\
                                    &=    \dis  \frac{4}{9}  \sen \theta + C\\
                                    &=    \dis  \frac{4}{9}  \frac{\sqrt{16x^2-9}}{4x} + C\\
                                    &=    \dis  \frac{\sqrt{16x^2-9}}{9x} + C\\
\end{align*}
\dotfill [1 punto]




\newpage
\item $\dis \int \sqrt{\e^{2t}-9} \, dt$\\

Utilizamos la siguiente sustitución trigonométrica:\\

\begin{minipage}{.3\textwidth}
\raggedright
\begin{align*}
e^t  &= 3\sec \theta      \\
e^t \, dt &= 3 \sec \theta \tan \theta  \, d\theta  \\
 dt &=  \tan \theta  \, d\theta  \\
\end{align*}
\end{minipage}% <---------------- Note the use of "%"
\begin{minipage}{.3\textwidth}
\raggedleft
\begin{tikzpicture}[thick]
\coordinate (O) at (0,0);
\coordinate (A) at (4,0);
\coordinate (B) at (0,2);
\draw (O)--(A)--(B)--cycle;

\tkzLabelSegment[below=2pt](O,A){$3$}
\tkzLabelSegment[left=2pt](O,B){$\sqrt{16x^2- 9}$}
\tkzLabelSegment[above right=2pt](A,B){$4x$}

\tkzMarkAngle[fill= orange,size=0.8cm,%
opacity=.4](B,A,O)
\tkzLabelAngle[pos = 0.6](B,A,O){$\theta$}

\end{tikzpicture}
\end{minipage}

\dotfill [1 punto]

Obtenemos lo siguiente:
\begin{align*}
 \dis \int \sqrt{\e^{2t}-9} \, dt &=    \dis \int \sqrt{9 \sec^2 \theta - 9} \tan \theta\, d\theta\\
 									&=    \dis \int 3\tan \theta \tan \theta\, d\theta\\
                                    &=    \dis \int 3\tan^2 \theta \, d\theta\\
                                    &=    \dis \int 3(\sec^2 \theta-1)  \, d\theta\\
                                    &=    \dis 3\int \sec^2 \theta  \, d\theta - 3\int  d\theta\\
                                    &=    \dis 3 \tan \theta - 3\theta + C\\
                                    &=    \dis 3 (\frac{\sqrt{e^{2t}-9}}{3}) - 3\arcsec (\frac{e^t}{2}) + C\\
\end{align*}

\dotfill [1 punto]

\end{enumerate}


\end{enumerate}

\vspace{10mm}




\end{document}
